\par\Needspace{6\baselineskip}\medskip
\begin{center}\hypertarget{native-222-article160}{}\pdfbookmark[2]{160.}{article-160}160.\end{center}
Bei der Darstellung des Verfahrens, dem die vorangegangenen Untersuchungen gewidmet waren, sind wir auf einige Sonderfälle gestoßen, in denen es sich zumindest in der von uns dargestellten Form nicht anwenden lässt. Diese Schwierigkeit tritt, wie wir gesehen haben, in den folgenden Fällen auf:

\SourceLeaf{223}\noindent\emph{erstens}, wenn einer der drei geozentrischen Orte entweder mit dem entsprechenden heliozentrischen Ort der Erde oder mit dem Gegenpunkt zusammenfällt (der letztere Fall kann offenbar nur eintreten, wenn der Himmelskörper zwischen Sonne und Erde hindurchgegangen ist),

\noindent\emph{zweitens}, wenn der erste geozentrische Ort des Himmelskörpers mit dem dritten zusammenfällt,

\noindent\emph{drittens}, wenn alle drei geozentrischen Orte zusammen mit dem zweiten heliozentrischen Ort der Erde auf demselben Großkreis liegen.

Im ersten Fall bleibt die Lage eines der Großkreise $AB$, $A'B'$, $A''B''$ unbestimmt, im zweiten und dritten Fall die Lage des Punktes $B^*$. Die oben dargestellten Verfahren, mit denen wir bei als bekannt vorausgesetzten Größen $P$, $Q$ aus den geozentrischen Orten die heliozentrischen zu bestimmen gelehrt haben, versagen daher in diesen Fällen. Dabei ist jedoch ein wesentlicher Unterschied zu beachten. Im ersten Fall ist diese Schwierigkeit allein dem Verfahren zuzuschreiben, im zweiten und dritten Fall liegt sie in der Natur der Aufgabe selbst. Im ersten Fall kann die Bestimmung also durchaus gelingen, wenn man das Verfahren geeignet abwandelt. Im zweiten und dritten Fall ist sie gänzlich unmöglich, und die heliozentrischen Orte bleiben unbestimmt. Es lohnt sich, diese Zusammenhänge kurz zu erläutern. Eine erschöpfende Behandlung wäre umso weniger zweckmäßig, als in all diesen Sonderfällen eine genaue Bestimmung der Bahn unmöglich ist, da sie schon durch die geringsten Beobachtungsfehler außerordentlich stark beeinflusst würde. Dieselbe Schwierigkeit tritt auch dann noch auf, wenn die Beobachtungen einem dieser Fälle zwar nicht genau entsprechen, ihm aber nahekommen. Darauf ist bei der Auswahl der Beobachtungen zu achten. Insbesondere darf kein Ort verwendet werden, an dem sich der Himmelskörper zugleich in der Nähe eines Knotens und in der Nähe der Opposition oder Konjunktion befindet. Ebenso sind Beobachtungen zu vermeiden, bei denen der Himmelskörper in der letzten Beobachtung fast zum selben geozentrischen Ort zurückgekehrt ist, den er bei der ersten einnahm, und schließlich solche, bei denen der vom mittleren heliozentrischen Ort der Erde zum mittleren geozentrischen Ort des Himmelskörpers gezogene Großkreis mit der Richtung der geozentrischen Bewegung einen sehr spitzen Winkel bildet und den ersten und dritten Ort gleichsam streift.
